\chapter*{مقدمة}
\addcontentsline{toc}{schapter}{المقدمة}
تُعد المصفوفات من أهم الأدوات الرياضية التي ظهرت في القرن التاسع عشر، حيث وفّرت وسيلة منظمة للتعامل مع الأنظمة الخطية والمعادلات متعددة المتغيرات. ومع تطور الرياضيات الحديثة، برزت الزمر كأحد المفاهيم الأساسية في الجبر التجريدي، إذ تُستخدم لدراسة البنى الرياضية التي تمتلك عمليات مغلقة وقوانين محددة مثل الانعكاس والتبديل. وعند الجمع بين هذين المفهومين، تظهر لنا زمرة المصفوفات التي تمثل مجموعة من المصفوفات تحت عملية الضرب، وتُعد مثالاً عملياً غنيّاً يوضح كيف يمكن للزمر أن تُطبّق في مجالات متعددة مثل الفيزياء النظرية، علم الحاسوب، والهندسة.  
يهدف هذا البحث إلى استعراض الأسس النظرية للمصفوفات والزمر، ثم الانتقال إلى دراسة زمرة المصفوفات بشكل خاص، مع بيان أهم خصائصها وتطبيقاتها، وذلك لتوضيح الدور المحوري الذي تلعبه هذه البنى الرياضية في تطوير العلوم الحديثة.